% RESULTADO_RECUPERADO: integral dimensional functor, eq:cZ-internos,
% eq:longitud-ruta, and definition of ell_0; composition with Article III.
% The consolidated proof and its algebraic check are specific to this contribution.
\subsection{Constitutive speed and kinematic realization of the clock}
\label{v:sec:enlace-velocidad-constitutiva}

Radiative propagation uses the speed section produced by the
angular transport of the same HMT state. The identification is established
before representation units are chosen. Let $x=A_{\mathrm{rad}}$ and
$y=C^*_{\mathrm{rad}}$ be the angular coordinates generated in
Section~\ref{v:ant:sec:ii-angulos}, with the channels
\eqref{v:ant:eq:ii-canales}, and let
$\mathsf R$ be the self-adjoint involution of the two sheets, with rank-one
projectors $P_+$ and $P_-$. On the domain $x>|y|$, the transport and its
channels are
\begin{equation}
 T=\exp(-xI-y\mathsf R)=q_+P_++q_-P_-,
 \qquad q_+=e^{-x-y},\quad q_-=e^{-x+y},\quad 0<q_\pm<1.
 \label{v:eq:velocidad-transporte}
\end{equation}
The response of the temporal blocks of lengths $90$ and $120$ is
\begin{equation}
 \mathsf V=(I-T^{90})(I-T^{120})^{-1}
       =r_+P_++r_-P_-,\qquad
 r_\pm=\frac{1-q_\pm^{90}}{1-q_\pm^{120}}.
 \label{v:eq:velocidad-respuesta}
\end{equation}
The construction starts from the inherited channels; the symbols $x,y$ and
$q_\pm$ denote their coordinates, not parameters fitted to a target
speed.

\begin{proposition}[Identity of the constitutive and kinematic realizations]
\label{v:prop:identidad-velocidad}
In the same dimensional chart, the clock section $c_{\mathrm{int}}$
and the constitutive-response section $c_{\mathrm{vac}}$ coincide:
\begin{equation}
 c_{\mathrm{int}}=c_{\mathrm{vac}}
   =\widehat c\,u_C,
 \qquad
 \widehat c=(r_+r_-)^{-1}
            =\exp(-\mathcal E),
 \label{v:eq:identidad-velocidad}
\end{equation}
where
\begin{equation}
 \mathcal E=
 \log\!\bigl[(1-q_+^{90})(1-q_-^{90})\bigr]
 -\log\!\bigl[(1-q_+^{120})(1-q_-^{120})\bigr].
 \label{v:eq:velocidad-lector-simetrico}
\end{equation}
The equality preserves sheet exchange.
\end{proposition}
\begin{proof}
The factors appearing in the logarithms are positive. Therefore,
\[
 \mathcal E=\log(r_+r_-)
            =\log\det\mathsf V,
 \qquad \exp(-\mathcal E)=(\det\mathsf V)^{-1}.
\]
The determinant is computed here in the two-sheet realization with rank-one
projectors. The kinematic definition $c_{\mathrm{int}}=\exp(-\mathcal E)u_C$
and the constitutive definition $c_{\mathrm{vac}}=(r_+r_-)^{-1}u_C$
therefore give the same section. Exchanging $P_+\leftrightarrow
P_-$ interchanges $r_+$ and $r_-$ and preserves their product. Finally,
$\det T=e^{-2x}$ corresponds to the elementary transport, whereas
$\det\mathsf V=r_+r_-$ corresponds to its temporal response; the two
determinants have distinct roles in this composition.
\end{proof}

\subsubsection{Trajectory length and elementary duration}

The positive dimensional frame $(u_S,u_C,u_T,u_Q)$ induces the length
basis $u_L=u_Cu_T$. For an enriched trajectory $\gamma$ of
$n=|\gamma|$ transitions with a constant channel, the additive readers are
\begin{equation}
 T_{\mathrm{int}}(\gamma)=n u_T,
 \qquad L_{\mathrm{int}}(\gamma)=n\widehat c\,u_L.
 \label{v:eq:velocidad-longitud-reloj}
\end{equation}
Their additivity concerns the count and realization of each transition:
enriched routes additionally preserve orientation and memory.

\begin{corollary}[Normalization of the elementary length]
\label{v:cor:longitud-elemental-constitutiva}
For $n>0$, $L_{\mathrm{int}}(\gamma)/T_{\mathrm{int}}(\gamma)
=c_{\mathrm{vac}}$. In the normalization $t_0=u_T$,
\begin{equation}
 \ell_0=c_{\mathrm{int}}t_0=\widehat c\,u_L.
 \label{v:eq:longitud-elemental-constitutiva}
\end{equation}
In a general temporal chart $t_0=\tau_0u_T$, with $\tau_0>0$,
the corresponding expression is $\ell_0=\tau_0\widehat c\,u_L$.
\end{corollary}
\begin{proof}
Dividing the expressions in \eqref{v:eq:velocidad-longitud-reloj} cancels the integer
$n$ and uses $u_L/u_T=u_C$. The elementary formula is obtained
by multiplying the speed section by $t_0$.
\end{proof}

Generalization to an edge-dependent channel preserves the additive
formula $L_{\mathrm{int}}(\gamma)=\sum_{e\in\gamma}\widehat c(e)u_L$.
Its quotient by $nu_T$ is the mean of the speeds of those edges.
The homogeneous radiative realization uses the constant channel specified
in \eqref{v:eq:velocidad-longitud-reloj}.

\subsubsection{Change of chart and adapted normalization}

\begin{proposition}[Covariance of the identification]
\label{v:prop:covariancia-identidad-velocidad}
Let $u_C'=d u_C$ and $u_T'=\eta u_T$, with $d,\eta>0$.
The same speed and the same elementary duration have coordinates
\begin{equation}
 \widehat c'=d^{-1}\widehat c,
 \qquad \tau_0'=\eta^{-1}\tau_0,
 \qquad u_L'=d\eta u_L.
 \label{v:eq:velocidad-cambio-carta}
\end{equation}
In particular, if $c_V=\widehat c_Vu_C^V$ and
$u_C^V=d u_C^{III}$, equality of the sections is exactly equivalent to
$d\widehat c_V=\widehat c_{III}$.
\end{proposition}
\begin{proof}
Express each section in the two bases. This gives
$\widehat c'u_C'=\widehat c u_C$ and
$\tau_0'u_T'=\tau_0u_T$. Their product satisfies
\[
 \tau_0'\widehat c' u_L'
 =\frac{\tau_0}{\eta}\frac{\widehat c}{d}
       d\eta u_L
 =\tau_0\widehat c u_L=\ell_0.
\]
The final equivalence follows by substituting the relation between the bases.
\end{proof}

The adapted choice $d=\widehat c$ gives $u_C'=c_{\mathrm{vac}}$ and
$\widehat c'=1$. This choice expresses the already constructed section in a
basis adapted to it. The spectral coefficient in the internal chart remains
$(r_+r_-)^{-1}$; spectral inversion uses that internal chart.
For example, retaining the action and charge bases, the constitutive
coordinates transform as
\[
 \widehat\mu'=d\widehat\mu,
 \qquad \widehat\varepsilon'=d\widehat\varepsilon,
 \qquad
 \widehat\mu'\widehat\varepsilon'=(\widehat c')^{-2}.
\]
In the adapted chart, $\widehat\mu'=r_-/r_+$ and
$\widehat\varepsilon'=r_+/r_-$, since the internal coordinates are
$\widehat\mu=r_-^2$ and $\widehat\varepsilon=r_+^2$.

\subsubsection{Oriented Catalan moment and functional bridge}

The constitutive response does not receive Catalan's constant as an
external number. The dodecaphasic quarter turn produces the collection of
oriented moments
\begin{equation}
 \mathcal C_2(s)=\sum_{k\geq0}
 \frac{(-1)^ks^{2k+1}}{(2k+1)^2},\qquad
 \mathcal D=s\frac{\dd}{\dd s},\qquad 0<s<1.
 \label{v:eq:momento-catalan}
\end{equation}
Absolute convergence on compact subsets permits termwise differentiation
and gives
\begin{equation}
 \mathcal D^2\mathcal C_2(s)=\frac{s}{1+s^2},\qquad
 f(s)=\frac{1+s+s^2}{s(1+s)}\mathcal D^2\mathcal C_2(s)
 =\frac{1-s^3}{1-s^4}.
 \label{v:eq:puente-catalan-constitutivo}
\end{equation}
With $s_\pm=q_\pm^{30}$, the incidences $90=3\cdot30$ and
$120=4\cdot30$ produce exactly
\begin{equation}
 r_\pm=f(s_\pm),\qquad
 \mathcal C_2(1)=G_{\mathrm{Cat}}.
 \label{v:eq:catalan-canales}
\end{equation}
The endpoint value is the representation of the collection; the functional
collection, rather than that value alone, governs the two channels.

The bridge is reversible within the fixed regularity class. Given $f$,
the condition $\mathcal C_2(s)\to0$ as $s\downarrow0$ uniquely selects
\begin{equation}
 \mathcal C_2(s)=\int_0^s\log\frac{s}{t}\,
 \frac{1+t}{1+t+t^2}f(t)\,\dd t.
 \label{v:eq:restitucion-funcional-catalan}
\end{equation}
Indeed, this integral satisfies
$\mathcal D^2\mathcal C_2=s(1+s)f/(1+s+s^2)$; for
\eqref{v:eq:puente-catalan-constitutivo}, the integrand reduces to
$\log(s/t)/(1+t^2)$. Two homogeneous solutions differ by
$a\log s+b$, and null regularity at the origin forces $a=b=0$.

The same dilogarithmic collection organizes a joint potential. Let
$\mathcal D_2(z)=\sum_{n\geq1}z^n/n^2$ and
\begin{equation}
 \mathscr F(x,y)=\sum_{\sigma=\pm}
 \left[
 \frac{\mathcal D_2(q_\sigma^{120})}{120}
 -\frac{\mathcal D_2(q_\sigma^{90})}{90}
 \right].
 \label{v:eq:potencial-constitutivo-catalan}
\end{equation}
Since $\dd\mathcal D_2(z)/\dd\log z=-\log(1-z)$ and
$q_\pm=e^{-x\mp y}$,
\begin{equation}
 \partial_x\mathscr F=\log\widehat c,
 \qquad
 \partial_y\mathscr F=\log\widehat Z,
 \qquad \widehat Z=\frac{r_-}{r_+}.
 \label{v:eq:gradiente-constitutivo-catalan}
\end{equation}
Thus speed and impedance are the two conjugate derivatives of a single
potential, whose additive constant is fixed by $\mathscr F\to0$ as
$q_\pm\to0$.

\subsubsection{Joint reconstruction from speed and Barbero--Immirzi}

Let $\psi=f^{-1}$ and define
\begin{equation}
 \mathcal H_B(s)=\frac{12s^3}{1-s^9}
 -\frac{s^4}{1-s^{12}}-\frac{(\log s)^2}{20\pi}.
 \label{v:eq:funcional-barbero-espectral}
\end{equation}
The formula in \eqref{v:ant:eq:gamma}, together with
\eqref{v:eq:catalan-canales}, factorizes exactly as
\begin{equation}
 \gamma_{\HMT}=\mathcal H_B(\psi(r_-))
 -\mathcal H_B(\psi(r_+)).
 \label{v:eq:factorizacion-barbero-espectral}
\end{equation}

\begin{theorem}[Joint reconstruction of the labelled channels]
\label{v:thm:restitucion-c-gamma}
In the normalized constitutive collection, the map
\begin{equation}
 \Psi:(3/4,1)^2\longrightarrow(1,16/9)\times\mathbb R,
 \qquad
 (r_+,r_-)\longmapsto(\widehat c,\gamma_{\HMT})
 \label{v:eq:mapa-c-gamma}
\end{equation}
is a diffeomorphism. Hence $(\widehat c,\gamma_{\HMT})$ reconstructs
the two labelled eigenvalues; the projectors $P_\pm$ remain data of the
initial representation.
\end{theorem}
\begin{proof}
The regular expression $f(s)=(1+s+s^2)/(1+s+s^2+s^3)$ satisfies
\[
 f'(s)=-\frac{s^2(s^2+2s+3)}{(1+s+s^2+s^3)^2}<0,
\]
with $f(0^+)=1$ and $f(1^-)=3/4$. Moreover,
\begin{equation}
 \mathcal H_B'(s)=
 \frac{36s^2(1+2s^9)}{(1-s^9)^2}
 -\frac{4s^3(1+2s^{12})}{(1-s^{12})^2}
 -\frac{\log s}{10\pi s}>0.
 \label{v:eq:monotonia-hb}
\end{equation}
The ratio of the first positive term to the second before subtraction is
\[
 \frac9s\frac{1+2s^9}{1+2s^{12}}
 \left(\frac{1-s^{12}}{1-s^9}\right)^2>9,
\]
and the last summand is also positive. Therefore
$F_B=\mathcal H_B\circ\psi$ is strictly decreasing, with limits
$F_B(3/4^+)=+\infty$ and $F_B(1^-)=-\infty$.

Fix $c_0=\widehat c$ and write $z=\log\widehat Z$. Then
$r_\pm=c_0^{-1/2}e^{\mp z/2}$ and
\[
 |z|<a(c_0):=\min\!\left\{\log c_0,
 \log\frac{16}{9c_0}\right\}.
\]
The function $\Gamma_{c_0}(z)=F_B(r_-)-F_B(r_+)$ satisfies
\[
 \Gamma_{c_0}'(z)=\frac12
 \bigl[r_-F_B'(r_-)+r_+F_B'(r_+)\bigr]<0
\]
and ranges over $\mathbb R$ between the two admissible endpoints. There is
a unique solution for every $\gamma_{\HMT}$; the nonzero derivative and
the implicit function theorem provide the smooth global inverse.
\end{proof}

On the image produced by the HMT generator, reconstruction preserves
causal order. Once $r_\pm$ have been recovered, one obtains
$s_\pm=\psi(r_\pm)$ and
\begin{equation}
 A_{\deg}=-\frac3\pi\log(s_+s_-),\qquad
 C^*_{\deg}=\frac3\pi\log\frac{s_-}{s_+},
 \label{v:eq:restitucion-angular-constitutiva}
\end{equation}
whence, on the already generated positive orientation,
$\alpha_{\HMT}=A_{\deg}/1000$ and
$\eta_{\mathrm{ret}}=C^*_{\deg}/2-\alpha_{\HMT}$. This path reconstructs
representations of the state; it does not turn
$(\widehat c,\gamma_{\HMT})$ or conventional values into inputs of the
generator.

\subsubsection{Transport to the radiative sector}

The dispersion relation of the homogeneous radiative realization is written
$\omega_{\symbf{k}}=c_{\mathrm{vac}}|\symbf{k}|$.
With $h=2\pi\hbar$ and $\beta=(k_B\Theta)^{-1}$, its spectral density
and Stefan--Boltzmann coefficient have the expressions
\begin{equation}
 u_\omega=
 \frac{\hbar\omega^3}{\pi^2(\widehat c u_C)^3}
 \frac{1}{e^{\beta\hbar\omega}-1},
 \qquad
 \sigma=\frac{\pi^2k_B^4}{60\hbar^3(\widehat c u_C)^2}.
 \label{v:eq:radiacion-velocidad-constitutiva}
\end{equation}
These substitutions transport precisely the same speed
section. The conditions of the radiative realization are retained:
a free three-dimensional bosonic sector, two transverse polarizations,
zero chemical potential, and positive temperature. Identity
\eqref{v:eq:identidad-velocidad} establishes its constitutive antecedent;
mode counting and the thermodynamic limit belong to the specific
proofs for that realization.

\paragraph{Normalization under amplification.}
If refinement acts by $\mathsf V_m=\mathsf V\otimes I_{9^m}$,
the normalized trace $\tau_m=\operatorname{Tr}/(2\cdot9^m)$ satisfies
\[
 \exp[-2\tau_m(\log\mathsf V_m)]
 =\exp[-\log r_+-\log r_-]=\widehat c.
\]
Indeed, each eigenvalue appears $9^m$ times. The multiplicity
cancels in the normalized trace. The ordinary amplified determinant is
$(r_+r_-)^{9^m}$: the two-sheet normalization, or equivalently the
normalized trace, preserves the constitutive reader under that
refinement. This amplification of the representation preserves the
enriched transport from which it arises.
